% GENERATED FILE. Edit canonical lessons, then run npm run generate:books.
% canonical-source-hash: fnv1a64:92079bfa045d902b
% canonical-lessons: RU-C441-prokhozhiy, RU-C441-rassprashivat, RU-C441-podskazyvat, RU-C441-razyskivat, RU-C441-peresprashivat

\chapter{Asking the way}
\label{ch:ru-asking-the-way}
\hlchaptermodality{441}

\begin{chapteropening}
\textbf{By the end of this chapter:} \emph{I can ask the way: a passer-by, asking around, telling someone the way, tracking a place down and asking again.}

Complete the last lesson of chapter 441: I can ask the way: a passer-by, asking around, telling someone the way, tracking a place down and asking again.
\end{chapteropening}

\section[\texorpdfstring{\ru{прохожий} (prokhózhiy)}{прохожий (prokhózhiy)}]{\ru{прохожий} (prokhózhiy) — a passer-by}

\label{lesson:RU-C441-prokhozhiy}



\begin{quote}
Before the new one: say the Russian for a boulevard, then the Russian for an avenue.
\end{quote}

\subsection*{You'll want to know: \ru{прохожий}}
\textbf{\ru{прохожий}} (\emph{prokhózhiy}) — \textquotedblleft{}a passer-by\textquotedblright{}.

\begin{quote}
\ru{Мы} \ru{спросили} \ru{дорогу} \ru{у} \ru{прохожего}.
\end{quote}

\begin{quote}
\emph{My sprosíli dorógu u prokhózhego.}
\end{quote}

\begin{quote}
We asked a passer-by the way.
\end{quote}

An adjective used as a noun: \emph{\ru{прохожего}}, \emph{\ru{прохожие}}.

\begin{cousinweb}
\emph{\ru{прохожий}} is built on \emph{\ru{проходить}} 'to go past' with an adjective ending: someone walking by. \emph{\ru{Ходить}} goes back to Proto-Slavic \emph{xodŭ} 'motion' (EtymDB).
\end{cousinweb}

\begin{grammarlens}[title={What you've built}]
One more word for telling, explaining and describing at B1.
\end{grammarlens}

\subsection*{Your turn}
\begin{itemize}
\raggedright
  \item \emph{Say it:} \emph{prokhózhiy}
  \item \emph{Say it:} \emph{prokhózhiy}, once more
  \item \emph{Say it:} a sentence of your own with \emph{prokhózhiy}
  \item \emph{Recall:} say the Russian for a boulevard, then the Russian for an avenue, then say \emph{prokhózhiy} again
\end{itemize}

\begin{tcolorbox}[breakable,colback=teal!4,colframe=teal!35!black,title={Before you move on}]
What does \ru{прохожий} mean? (\textquotedblleft{}A passer-by\textquotedblright{}.) Say it once more.
\end{tcolorbox}
\section[\texorpdfstring{\ru{расспрашивать} (rasspráshivat)}{расспрашивать (rasspráshivat)}]{\ru{расспрашивать} (rasspráshivat) — to ask around, to question}

\label{lesson:RU-C441-rassprashivat}



\begin{quote}
Before the new one: say the Russian for an avenue, then the Russian for a passer-by.
\end{quote}

\subsection*{You'll want to know: \ru{расспрашивать}}
\textbf{\ru{расспрашивать}} (\emph{rasspráshivat}) — \textquotedblleft{}to ask around, to question\textquotedblright{}.

\begin{quote}
\ru{Туристы} \ru{расспрашивают} \ru{местных} \ru{о} \ru{дороге}.
\end{quote}

\begin{quote}
\emph{Turísty rasspráshivayut méstnykh o doróge.}
\end{quote}

\begin{quote}
The tourists ask the locals about the way.
\end{quote}

Imperfective; perfective \emph{\ru{расспросить}}. Takes the accusative of the person and \emph{\ru{о}} + prepositional.

\begin{cousinweb}
\emph{\ru{расспрашивать}} is built from \emph{\ru{рас}-} 'all round' and \emph{\ru{спрашивать}} 'to ask', on the root of \emph{\ru{просить}} 'to ask', which EtymDB takes from Proto-Slavic \emph{prositi} and the PIE root \emph{pre\'{k}-}.
\end{cousinweb}

\begin{grammarlens}[title={What you've built}]
One more word for telling, explaining and describing at B1.
\end{grammarlens}

\subsection*{Your turn}
\begin{itemize}
\raggedright
  \item \emph{Say it:} \emph{rasspráshivat}
  \item \emph{Say it:} \emph{rasspráshivat}, once more
  \item \emph{Say it:} a sentence of your own with \emph{rasspráshivat}
  \item \emph{Recall:} say the Russian for an avenue, then the Russian for a passer-by, then say \emph{rasspráshivat} again
\end{itemize}

\begin{tcolorbox}[breakable,colback=teal!4,colframe=teal!35!black,title={Before you move on}]
What does \ru{расспрашивать} mean? (\textquotedblleft{}To ask around, to question\textquotedblright{}.) Say it once more.
\end{tcolorbox}
\section[\texorpdfstring{\ru{подсказывать} (podskázyvat)}{подсказывать (podskázyvat)}]{\ru{подсказывать} (podskázyvat) — to tell (someone the way); to prompt}

\label{lesson:RU-C441-podskazyvat}



\begin{quote}
Before the new one: say the Russian for a passer-by, then the Russian for to ask around.
\end{quote}

\subsection*{You'll want to know: \ru{подсказывать}}
\textbf{\ru{подсказывать}} (\emph{podskázyvat}) — \textquotedblleft{}to tell (someone the way); to prompt\textquotedblright{}.

\begin{quote}
\ru{Прохожий} \ru{подсказывает} \ru{нам} \ru{дорогу} \ru{к} \ru{вокзалу}.
\end{quote}

\begin{quote}
\emph{Prokhózhiy podskázyvayet nam dorógu k vokzálu.}
\end{quote}

\begin{quote}
A passer-by tells us the way to the station.
\end{quote}

Imperfective; perfective \emph{\ru{подсказать}} (\emph{\ru{подскажу}, \ru{подскажешь}}). Takes the dative: \emph{\ru{подскажите}, \ru{пожалуйста}} — could you tell me, please.

\begin{cousinweb}
\emph{\ru{подсказывать}} is built from \emph{\ru{под}-} 'under, quietly' and \emph{-\ru{сказывать}}, from \emph{\ru{сказать}} 'to say': to say something quietly to help. EtymDB lists \emph{\ru{сказать}} without a deeper source.
\end{cousinweb}

\begin{grammarlens}[title={What you've built}]
One more word for telling, explaining and describing at B1.
\end{grammarlens}

\subsection*{Your turn}
\begin{itemize}
\raggedright
  \item \emph{Say it:} \emph{podskázyvat}
  \item \emph{Say it:} \emph{podskázyvat}, once more
  \item \emph{Say it:} a sentence of your own with \emph{podskázyvat}
  \item \emph{Recall:} say the Russian for a passer-by, then the Russian for to ask around, then say \emph{podskázyvat} again
\end{itemize}

\begin{tcolorbox}[breakable,colback=teal!4,colframe=teal!35!black,title={Before you move on}]
What does \ru{подсказывать} mean? (\textquotedblleft{}To tell (someone the way); to prompt\textquotedblright{}.) Say it once more.
\end{tcolorbox}
\section[\texorpdfstring{\ru{разыскивать} (razýskivat)}{разыскивать (razýskivat)}]{\ru{разыскивать} (razýskivat) — to search for, to track down}

\label{lesson:RU-C441-razyskivat}



\begin{quote}
Before the new one: say the Russian for to ask around, then the Russian for to tell.
\end{quote}

\subsection*{You'll want to know: \ru{разыскивать}}
\textbf{\ru{разыскивать}} (\emph{razýskivat}) — \textquotedblleft{}to search for, to track down\textquotedblright{}.

\begin{quote}
\ru{Мы} \ru{долго} \ru{разыскивали} \ru{эту} \ru{улицу}.
\end{quote}

\begin{quote}
\emph{My dólgo razýskivali étu úlitsu.}
\end{quote}

\begin{quote}
We searched for this street for a long time.
\end{quote}

Imperfective; perfective \emph{\ru{разыскать}} (\emph{\ru{разыщу}, \ru{разыщешь}}).

\begin{cousinweb}
\emph{\ru{разыскивать}} is built from \emph{\ru{раз}-} and \emph{-\ru{ыскивать}}, from \emph{\ru{искать}} 'to seek', which is inherited from Proto-Slavic \emph{jĭskati} (EtymDB).
\end{cousinweb}

\begin{grammarlens}[title={What you've built}]
One more word for telling, explaining and describing at B1.
\end{grammarlens}

\subsection*{Your turn}
\begin{itemize}
\raggedright
  \item \emph{Say it:} \emph{razýskivat}
  \item \emph{Say it:} \emph{razýskivat}, once more
  \item \emph{Say it:} a sentence of your own with \emph{razýskivat}
  \item \emph{Recall:} say the Russian for to ask around, then the Russian for to tell, then say \emph{razýskivat} again
\end{itemize}

\begin{tcolorbox}[breakable,colback=teal!4,colframe=teal!35!black,title={Before you move on}]
What does \ru{разыскивать} mean? (\textquotedblleft{}To search for, to track down\textquotedblright{}.) Say it once more.
\end{tcolorbox}
\section[\texorpdfstring{\ru{переспрашивать} (perespráshivat)}{переспрашивать (perespráshivat)}]{\ru{переспрашивать} (perespráshivat) — to ask again, to ask (someone) to repeat}

\label{lesson:RU-C441-peresprashivat}



\begin{quote}
Before the new one: say the Russian for to tell, then the Russian for to search for.
\end{quote}

\subsection*{You'll want to know: \ru{переспрашивать}}
\textbf{\ru{переспрашивать}} (\emph{perespráshivat}) — \textquotedblleft{}to ask again, to ask (someone) to repeat\textquotedblright{}.

\begin{quote}
\ru{Если} \ru{я} \ru{не} \ru{понимаю} \ru{ответ}, \ru{я} \ru{переспрашиваю}.
\end{quote}

\begin{quote}
\emph{Yésli ya ne ponimáyu otvét, ya perespráshivayu.}
\end{quote}

\begin{quote}
If I do not understand the answer, I ask again.
\end{quote}

Imperfective; the perfective is \emph{\ru{переспросить}}, with \emph{\ru{переспрошу}, \ru{переспросишь}}. It takes the accusative of the person: \emph{\ru{переспросить} \ru{прохожего}}.

\begin{cousinweb}
\emph{\ru{переспрашивать}} is built from \emph{\ru{пере}-} 'again' and \emph{\ru{спрашивать}} 'to ask', on the root of \emph{\ru{просить}} 'to ask', which EtymDB takes from Proto-Slavic \emph{prositi} and the PIE root \emph{pre\'{k}-}.
\end{cousinweb}

\begin{grammarlens}[title={What you've built}]
One more word for telling, explaining and describing at B1.
\end{grammarlens}

\subsection*{Your turn}
\begin{itemize}
\raggedright
  \item \emph{Say it:} \emph{perespráshivat}
  \item \emph{Say it:} \emph{perespráshivat}, once more
  \item \emph{Say it:} a sentence of your own with \emph{perespráshivat}
  \item \emph{Recall:} say the Russian for to tell, then the Russian for to search for, then say \emph{perespráshivat} again
\end{itemize}

\begin{tcolorbox}[breakable,colback=teal!4,colframe=teal!35!black,title={Before you move on}]
What does \ru{переспрашивать} mean? (\textquotedblleft{}To ask again, to ask (someone) to repeat\textquotedblright{}.) Say it once more.
\end{tcolorbox}
